%% toc: P15 174 174.1,174.2,174.3,174.4
\section{Proofs for target populations, sampling and observational comparisons}
\label{sec:annex-p15-w02b-p15-075}

Notation is that of
\cref{def:p15-w02b-p15-075-populations,def:p15-w02b-p15-075-complete,def:p15-w02b-p15-075-sampling-design}.
Two elementary facts are used repeatedly. For real numbers \(x_1,\ldots,x_n\)
with mean \(\bar x\) and any real \(c\),
\[
 \sum_i(x_i-c)^2=\sum_i(x_i-\bar x)^2+n(\bar x-c)^2, \tag{A}
\]
because the cross term \(2(\bar x-c)\sum_i(x_i-\bar x)\) vanishes. For a real
random variable \(X\) with finite variance and \(a>0\), Markov's inequality
applied to \((X-\mathbb EX)^2\) gives
\[
 \Pr(|X-\mathbb EX|>a)\leq\Pr\bigl((X-\mathbb EX)^2\geq a^2\bigr)
 \leq\operatorname{Var}(X)/a^2, \tag{B}
\]
since \(a^2\mathbf1_{\{(X-\mathbb EX)^2\geq a^2\}}\leq(X-\mathbb EX)^2\)
pointwise.

\begin{annexitem}{annex:p15-w02b-p15-075-estimand-gap}{thm:p15-w02b-p15-075-estimand-gap}{Composition, transport and the target range}
\begin{proof}
(1) Grouping units by their covariate value gives
\[
 \tau=\sum_{v:\,q_{\mathcal R}(v)>0}q_{\mathcal R}(v)\tau(v),\qquad
 \tau_{\mathcal U}=\sum_{v:\,q_{\mathcal U}(v)>0}q_{\mathcal U}(v)\tau_{\mathcal U}(v).
\]
By hypothesis every \(v\) with \(q_{\mathcal U}(v)>0\) has
\(q_{\mathcal R}(v)>0\), so \(\tau(v)\) is defined there, and the number
\[
 A=\sum_{v:\,q_{\mathcal U}(v)>0}q_{\mathcal U}(v)\tau(v)
  =\sum_{v:\,q_{\mathcal R}(v)>0}q_{\mathcal U}(v)\tau(v)
\]
is defined; the two sums agree because terms with \(q_{\mathcal U}(v)=0\)
contribute nothing. Then
\(\tau_{\mathcal U}-\tau=(A-\tau)+(\tau_{\mathcal U}-A)\). The first bracket
is the composition sum and the second is the within-stratum sum.

(2) Under the stated equalities the within-stratum sum vanishes, so
\(\tau_{\mathcal U}=A\), which is the displayed formula. If also
\(q_{\mathcal U}=q_{\mathcal R}\), the composition sum vanishes and
\(\tau_{\mathcal U}=\tau\).

(3) Split the sum defining \(\tau_{\mathcal U}\) into the \(N\) recruited and
the \(M-N\) unrecruited units:
\(M\tau_{\mathcal U}=N\tau+(M-N)\tau_{\mathcal U\setminus\mathcal R}\).
Division by \(M\) gives the identity. If each unrecruited contrast lies in
\([c_-,c_+]\), so does their average, which gives the two endpoints. Every
value \(\theta\in[c_-,c_+]\) is attained as that average by setting every
unrecruited contrast equal to \(\theta\). Hence every value of the displayed
interval is compatible, and no other value is.

(4) If the sets are disjoint, the recruited units' potential responses do not
enter \(\tau_{\mathcal U}\). Setting every target contrast equal to
\(\theta\in[c_-,c_+]\) gives \(\tau_{\mathcal U}=\theta\).
\end{proof}
\end{annexitem}

\begin{annexitem}{annex:p15-w02b-p15-075-neyman}{thm:p15-w02b-p15-075-neyman}{The variance of a completely randomized contrast}
\begin{proof}
Write \(A_i(t)=\mathbf1_{\{Z_i=t\}}\), \(p_t=n_t/N\) and
\(e_i(t)=Y_i(t)-\bar y(t)\), so \(\sum_ie_i(t)=0\).

\emph{Assignment moments.} By symmetry of the uniform law, \(\Pr(Z_i=t)=p_t\).
For \(i\neq j\), \(\Pr(Z_i=t,Z_j=t)=n_t(n_t-1)/(N(N-1))\) and, for
\(t\neq u\), \(\Pr(Z_i=t,Z_j=u)=n_tn_u/(N(N-1))\). Subtracting products
of marginals gives
\[
\begin{aligned}
 \operatorname{Cov}(A_i(t),A_i(t))&=p_t(1-p_t), &
 \operatorname{Cov}(A_i(t),A_i(u))&=-p_tp_u,\\
 \operatorname{Cov}(A_i(t),A_j(t))&=-\frac{p_t(1-p_t)}{N-1}, &
 \operatorname{Cov}(A_i(t),A_j(u))&=\frac{p_tp_u}{N-1},
\end{aligned}
\]
for \(i\neq j\) and \(t\neq u\). For instance,
\(n_t(n_t-1)/(N(N-1))-p_t^2=p_t\bigl((n_t-1)/(N-1)-n_t/N\bigr)
=-p_t(1-p_t)/(N-1)\).

(1) Since \(\overline Y_t=n_t^{-1}\sum_iA_i(t)Y_i(t)\) and
\(\mathbb EA_i(t)=p_t\), \(\mathbb E\overline Y_t=\bar y(t)\), and so
\(\mathbb E\widehat\tau=\sum_tw_t\bar y(t)=\tau\). Moreover
\(\sum_iA_i(t)=n_t\), hence
\[
 \widehat\tau-\tau=\sum_t\frac{w_t}{n_t}\sum_iA_i(t)e_i(t),
\]
and
\[
 \operatorname{Var}(\widehat\tau)=\sum_{t,u}\frac{w_tw_u}{n_tn_u}
 \sum_{i,j}e_i(t)e_j(u)\operatorname{Cov}(A_i(t),A_j(u)).
\]
For fixed \(t,u\), use \(\sum_{i\neq j}e_i(t)e_j(u)
=\bigl(\sum_ie_i(t)\bigr)\bigl(\sum_je_j(u)\bigr)-\sum_ie_i(t)e_i(u)
=-\sum_ie_i(t)e_i(u)\). Put
\(S_{tu}=(N-1)^{-1}\sum_ie_i(t)e_i(u)\), so \(S_{tt}=S_t^2\). For
\(t=u\) the inner double sum equals
\(p_t(1-p_t)\bigl(1+\tfrac1{N-1}\bigr)\sum_ie_i(t)^2=Np_t(1-p_t)S_t^2\).
For \(t\neq u\) it equals
\(-p_tp_u\bigl(1+\tfrac1{N-1}\bigr)\sum_ie_i(t)e_i(u)=-Np_tp_uS_{tu}\).
Multiplying by \(w_tw_u/(n_tn_u)\) and using \(p_t=n_t/N\) gives
\(w_t^2S_t^2(1/n_t-1/N)\) for \(t=u\) and \(-w_tw_uS_{tu}/N\) for
\(t\neq u\). Hence
\[
 \operatorname{Var}(\widehat\tau)=\sum_t\frac{w_t^2S_t^2}{n_t}
 -\frac1N\sum_{t,u}w_tw_uS_{tu}.
\]
Finally \(\tau_i-\tau=\sum_tw_te_i(t)\), so
\[
 S_w^2=\frac1{N-1}\sum_i\Bigl(\sum_tw_te_i(t)\Bigr)^2=\sum_{t,u}w_tw_uS_{tu}.
\]

(2) The set \(\{i:Z_i=t\}\) is a uniformly random subset of size \(n_t\)
of \(\{1,\ldots,N\}\), and on it the observed responses are the fixed values
\(Y_i(t)\). Apply (A) with \(c=\bar y(t)\) to the \(n_t\) sampled values:
\[
 (n_t-1)s_t^2=\sum_{i:Z_i=t}e_i(t)^2-n_t\bigl(\overline Y_t-\bar y(t)\bigr)^2.
\]
Each unit is included with probability \(p_t\), so the first sum has
expectation \(p_t(N-1)S_t^2\). The second term has expectation
\(n_t\operatorname{Var}(\overline Y_t)\). By part~(1) applied with a single
nonzero weight \(w_t=1\), this variance is \(S_t^2(1/n_t-1/N)\), because the
unit contrasts are then \(Y_i(t)\) and \(S_w^2=S_t^2\). Thus
\[
 (n_t-1)\,\mathbb E[s_t^2]
 =S_t^2\Bigl(\frac{n_t(N-1)}N-1+\frac{n_t}N\Bigr)=(n_t-1)S_t^2,
\]
and \(n_t\geq2\) gives \(\mathbb E[s_t^2]=S_t^2\). Linearity and part~(1)
give \(\mathbb E\widehat V=\sum_tw_t^2S_t^2/n_t
=\operatorname{Var}(\widehat\tau)+S_w^2/N\). The excess is zero exactly when
\(\sum_i(\tau_i-\tau)^2=0\), that is, when every \(\tau_i=\tau\).

(3) Fix levels \(t\neq u\) with \(w_tw_u\neq0\). For reals \(a,b\), let
\(\mathcal Y(a,b)\) be the table with \(Y_1(t)=a\), \(Y_1(u)=b\) and every
other potential response zero. Let \(g\) be any function of the observed data.
For an assignment \(z\), the observed vector is \((a,0,\ldots,0)\) if
\(z_1=t\), \((b,0,\ldots,0)\) if \(z_1=u\), and zero otherwise. Hence
\[
 \mathbb E_{\mathcal Y(a,b)}[g]=F_t(a)+F_u(b)+F_0,
\]
where \(F_t(a)=\sum_{z:z_1=t}\Pr(z)\,g(z,(a,0,\ldots,0))\), \(F_u\) is
defined likewise, and \(F_0\) collects the remaining assignments. Therefore
\[
 \mathbb E_{\mathcal Y(1,1)}[g]-\mathbb E_{\mathcal Y(1,0)}[g]
 -\mathbb E_{\mathcal Y(0,1)}[g]+\mathbb E_{\mathcal Y(0,0)}[g]=0 .
\]
For the variance, \(S_t^2\) depends on \(a\) alone, \(S_u^2\) on \(b\)
alone, and the other \(S_s^2\) vanish. The unit contrasts are
\(\tau_1=w_ta+w_ub\) and \(\tau_i=0\) for \(i\geq2\), with
\(\tau=\tau_1/N\). Then
\(\sum_i(\tau_i-\tau)^2=\tau_1^2\bigl((N-1)^2+(N-1)\bigr)/N^2=\tau_1^2(N-1)/N\),
so \(S_w^2=\tau_1^2/N\). The same second difference applied to
\(\operatorname{Var}(\widehat\tau)\) therefore equals
\(-N^{-2}\bigl((w_t+w_u)^2-w_t^2-w_u^2\bigr)=-2w_tw_u/N^2\neq0\). So the
expectation of \(g\) cannot equal \(\operatorname{Var}(\widehat\tau)\) at all
four tables.
\end{proof}
\end{annexitem}

\begin{annexitem}{annex:p15-w02b-p15-075-chebyshev}{thm:p15-w02b-p15-075-chebyshev}{The Chebyshev interval and a variance bound from the response range}
\begin{proof}
(1) By \cref{thm:p15-w02b-p15-075-neyman}(1), \(\mathbb E\widehat\tau=\tau\).
Apply (B) with \(a=\sqrt{B/\eta}>0\) when \(B>0\):
\(\Pr(|\widehat\tau-\tau|>a)\leq\operatorname{Var}(\widehat\tau)/a^2
\leq\eta\). If \(B=0\), then \(\widehat\tau=\tau\) almost surely and the
event has probability one.

(2) For fixed \(t\), apply (A) with \(c=(m_-+m_+)/2\):
\(\sum_i(Y_i(t)-\bar y(t))^2\leq\sum_i(Y_i(t)-c)^2\leq Nh^2/4\), since each
\(|Y_i(t)-c|\leq h/2\). Hence \(S_t^2\leq\frac N{N-1}\frac{h^2}4\). By
\cref{thm:p15-w02b-p15-075-neyman}(1),
\(\operatorname{Var}(\widehat\tau)\leq\sum_tw_t^2S_t^2/n_t\leq B\).
\end{proof}
\end{annexitem}

\begin{annexitem}{annex:p15-w02b-p15-075-randomization-test}{thm:p15-w02b-p15-075-randomization-test}{The randomization p-value is exact}
\begin{proof}
\emph{A lemma.} Let \(T\) be a random variable with finitely many values and
\(G(x)=\Pr(T\geq x)\). Then \(\Pr(G(T)\leq\alpha)\leq\alpha\). Indeed, list
the values \(v_1<\cdots<v_m\). The function \(j\mapsto G(v_j)\) is
nonincreasing, so \(J=\{j:G(v_j)\leq\alpha\}\) is empty or equals
\(\{j_0,\ldots,m\}\) for some \(j_0\). In the first case the probability is
zero. In the second, \(\{G(T)\leq\alpha\}=\{T\geq v_{j_0}\}\), whose
probability is \(G(v_{j_0})\leq\alpha\).

(1) Let \(\mathcal Y\) be the true table. If \(H\) is true, the table that
\(H\) computes from the realized \((z,y)\) is \(\mathcal Y\), whatever \(z\)
was realized, so \(y^H(z')=(Y_i(z'_i))_i\) for every \(z'\). Put
\(T=T_H(Z,(Y_i(Z_i))_i)\) and \(T^*=T_H(Z',(Y_i(Z'_i))_i)\). Since \(Z\)
and \(Z'\) have the same law, so do \(T\) and \(T^*\). The p-value is
\(G(T)\), with \(G(x)=\Pr(T^*\geq x)=\Pr(T\geq x)\), a function fixed by
\(\mathcal Y\) and the design. The assignment law has finite support, so
\(T\) takes finitely many values, and the lemma applies.

(2) Under the constant-effect model, \(H_{\delta_0}\) is true. By (1),
\(\Pr(p_{H_{\delta_0}}\leq\alpha)\leq\alpha\), and
\(\{\delta_0\notin C\}=\{p_{H_{\delta_0}}\leq\alpha\}\). Finally,
\(\tau=N^{-1}\sum_i(Y_i(1)-Y_i(0))=\delta_0\).
\end{proof}
\end{annexitem}

\begin{annexitem}{annex:p15-w02b-p15-075-blocked}{thm:p15-w02b-p15-075-blocked}{Blocked assignment as independent randomizations}
\begin{proof}
(1) Summing the unit contrasts block by block gives
\(N\tau=\sum_bN_b\tau_b\). Inside block \(b\) the assignment is a complete
randomization of \(N_b\geq|\mathcal T|\geq2\) units with cell sizes
\(n_{tb}\geq1\). By \cref{thm:p15-w02b-p15-075-neyman}(1),
\(\widehat\tau_b=\sum_tw_t\overline Y_{tb}\) has expectation \(\tau_b\) and
variance \(V_b\). Linearity gives \(\mathbb E\widehat\tau_{\mathcal B}=\tau\).

(2) The \(\widehat\tau_b\) are functions of independent assignments, so they
are independent, and the variance of the weighted sum is
\(\sum_b(N_b/N)^2V_b\).

(3) By \cref{thm:p15-w02b-p15-075-neyman}(2) applied in each block,
\(\mathbb E\sum_{t:\,w_t\neq0}w_t^2s_{tb}^2/n_{tb}=V_b+S_{wb}^2/N_b\).
Multiplying by \((N_b/N)^2\) and summing proves the claim.
\end{proof}
\end{annexitem}

\begin{annexitem}{annex:p15-w02b-p15-075-inclusion}{thm:p15-w02b-p15-075-inclusion}{Inverse-inclusion estimation}
\begin{proof}
Write \(I_k=\mathbf1_{\{k\in\mathcal S\}}\), so \(\mathbb EI_k=\rho_k\) and
\(\mathbb E[I_kI_l]=\rho_{kl}\).

(1) Since \(\widehat y=M^{-1}\sum_kI_ky_k/\rho_k\), its expectation is
\(M^{-1}\sum_ky_k\). Its variance is
\[
 \frac1{M^2}\sum_{k,l}\operatorname{Cov}(I_k,I_l)\frac{y_ky_l}{\rho_k\rho_l},
\]
and \(\operatorname{Cov}(I_k,I_l)=\rho_{kl}-\rho_k\rho_l\), also for \(k=l\).

(2) \(\widehat v=M^{-2}\sum_{k,l}I_kI_l\,a_{kl}\) with
\(a_{kl}=(\rho_{kl}-\rho_k\rho_l)y_ky_l/(\rho_{kl}\rho_k\rho_l)\). Taking
expectations replaces \(I_kI_l\) by \(\rho_{kl}\), which cancels the
denominator \(\rho_{kl}\) and returns the variance formula of (1).

(3) Each unit lies in exactly \(\binom{M-1}{n-1}\) of the
\(\binom Mn\) subsets of size \(n\), so \(\rho_k=n/M\) and
\(\widehat y=n^{-1}\sum_{k\in\mathcal S}y_k\). For \(M\geq2\) and
\(n\leq M\), take \(N=M\), two cells \(\{\mathrm{in},\mathrm{out}\}\) with
sizes \(n\) and \(M-n\), and weights \(w_{\rm in}=1\), \(w_{\rm out}=0\).
The labelled partition into sampled and unsampled units is then a complete
randomization, provided \(n<M\). By
\cref{thm:p15-w02b-p15-075-neyman}(1), with \(Y_k(\mathrm{in})=y_k\),
\(S_w^2=S_{\mathcal U}^2\) and the variance is
\(S_{\mathcal U}^2/n-S_{\mathcal U}^2/M\). If \(n=M\), the estimator is the
constant \(\bar y_{\mathcal U}\), and the formula gives zero. The result agrees
with \cref{thm:p15-without-replacement-convex}, whose \(\sigma_C^2\) equals
\((M-1)S_{\mathcal U}^2/M\).

(4) If \(\rho_k=0\), then \(k\notin\mathcal S\) almost surely, so the law of
\((\mathcal S,(y_l)_{l\in\mathcal S})\) is the same for two value vectors that
differ only at \(k\). Any function of the sample then has the same expectation
under both vectors, whereas their population means differ by \(1/M\) times the
change at \(k\).
\end{proof}
\end{annexitem}

\begin{annexitem}{annex:p15-w02b-p15-075-two-stage}{thm:p15-w02b-p15-075-two-stage}{Sampling then assignment}
\begin{proof}
(1) For \(k\in\mathcal U\) and \(t\in\mathcal T\),
\[
 \mathbb E\Bigl[\frac{\mathbf1_{\{k\in\mathcal S\}}\mathbf1_{\{Z_k=t\}}}
 {\rho_k\,\pi_k(t\mid\mathcal S)}\Bigr]
 =\sum_{s\ni k}p(s)\frac{\pi_k(t\mid s)}{\rho_k\,\pi_k(t\mid s)}
 =\frac{\rho_k}{\rho_k}=1 .
\]
The potential responses are fixed, and the observed \(Y_k\) equals \(Y_k(t)\)
on the event \(Z_k=t\). So the expectation of
\(\widehat\tau_{\mathcal U}\) is \(M^{-1}\sum_k\sum_tw_tY_k(t)=\tau_{\mathcal U}\).

(2) Here \(\rho_k=n/M\) and \(\pi_k(t\mid s)=n_t/n\). Their product is
\(n_t/M\), so \(\widehat\tau_{\mathcal U}=\sum_tw_tn_t^{-1}\sum_{k:Z_k=t}Y_k\).
Let \(A_t=\{k:Z_k=t\}\) and \(A_\circ=\mathcal U\setminus\mathcal S\). A
specified labelled partition \((A_t)_t,A_\circ\) with the stated sizes arises
when \(\mathcal S=\bigcup_tA_t\), which has probability \(1/\binom Mn\), and
the within-sample assignment is that partition, which has conditional probability
\(\prod_tn_t!/n!\). The product is \(\prod_tn_t!\,(M-n)!/M!\), the same for
every such partition. Hence \((A_t)_t,A_\circ\) is a complete randomization of
the \(M\) units of \(\mathcal U\) into the cells \(\mathcal T\cup\{\circ\}\),
with sizes \(n_t\) and \(M-n\geq1\). Give the extra cell weight zero and
potential response zero. Unit contrasts are unchanged, and
\cref{thm:p15-w02b-p15-075-neyman}(1), applied to these \(M\) units, gives the
variance. The extra cell contributes \(0\cdot S_\circ^2/(M-n)=0\). Part~(2) of
the same theorem requires \(n_t\geq2\) only for nonzero weights; it gives the
expectation of \(\widehat V\).
\end{proof}
\end{annexitem}

\begin{annexitem}{annex:p15-w02b-p15-075-licenses}{thm:p15-w02b-p15-075-licenses}{What sampling and assignment each license}
\begin{proof}
(1) Unbiasedness is \cref{thm:p15-w02b-p15-075-neyman}(1). The compatible
range is \cref{thm:p15-w02b-p15-075-estimand-gap}(3).

(2) With \(Z_k\) fixed, \(y_k=Y_k(Z_k)\) is a fixed value for each target
unit, and \cref{thm:p15-w02b-p15-075-inclusion}(1) gives unbiasedness for
\(M^{-1}\sum_ky_k\). Treat \(\mathcal U\) as the recruited set of an
observational comparison with assignment \((Z_k)_k\). Both levels occur, so
\cref{thm:p15-w02b-p15-075-observational}(1) applies and writes the population
difference of observed means as \(\tau_1+\sigma\) over \(\mathcal U\).
Moreover \(\tau_1-\tau_{\mathcal U}\) and \(\sigma\) are fixed numbers
determined by the table and \((Z_k)_k\). The sampling design does not enter
them, and they are nonzero in general. For instance, with the six units of
\cref{def:p15-w02b-p15-075-table} as \(\mathcal U\) and
\(Z_k=1\) exactly for \(k\leq3\),
\cref{thm:p15-w02b-p15-075-worked}(6) gives the population difference
\(-5/3\neq3/2=\tau_{\mathcal U}\).
\end{proof}
\end{annexitem}

\begin{annexitem}{annex:p15-w02b-p15-075-observational}{thm:p15-w02b-p15-075-observational}{Decomposition and logical range of the naive contrast}
\begin{proof}
(1) The treated units' observed responses are \(Y_i(1)\) and the untreated
units' are \(Y_i(0)\). Add and subtract \(n_1^{-1}\sum_{i:z_i=1}Y_i(0)\) in
\(D(z)=n_1^{-1}\sum_{i:z_i=1}Y_i(1)-n_0^{-1}\sum_{i:z_i=0}Y_i(0)\).

(2) If every contrast equals \(c\), then \(\tau_1(z)=c=\tau\) for every \(z\),
so \(D(z)-\tau=\sigma(z)\). Order the values as
\(a_{(1)}\leq\cdots\leq a_{(N)}\). For the \(n_1\) smallest values,
\(n_1^{-1}\sum_{j\leq n_1}a_{(j)}\leq a_{(n_1)}\leq a_{(n_1+1)}\leq
n_0^{-1}\sum_{j>n_1}a_{(j)}\). If equality held throughout, every
\(a_{(j)}\) with \(j\leq n_1\) would equal \(a_{(n_1)}\), every one with
\(j>n_1\) would equal \(a_{(n_1+1)}\), and these two values would be equal, so
all values would be equal. Since they are not, \(\sigma(z)<0\). The
largest-value case is symmetric. A tie-breaking rule fixed in advance makes the
selected \(z\) a single assignment, chosen with probability one.

(3) Write \(N\tau=\sum_iY_i(1)-\sum_iY_i(0)\). The observed data fix \(Y_i(1)=y_i\)
for treated units and \(Y_i(0)=y_i\) for untreated units. The remaining
\(n_0\) values \(Y_i(1)\) and \(n_1\) values \(Y_i(0)\) are unconstrained in
\([m_-,m_+]\). The expression is increasing in each unknown \(Y_i(1)\) and
decreasing in each unknown \(Y_i(0)\). Its minimum, attained at
\(Y_i(1)=m_-\) and \(Y_i(0)=m_+\), is \(NL(z)\); its maximum is \(NU(z)\).
The unknowns range over a box, which is connected, and \(\tau\) is continuous
in them, so every intermediate value is attained. The width is
\((n_0+n_1)h/N=h\). None of this used the assignment law, so the same range
holds for a randomized realization with the same observed data.
\end{proof}
\end{annexitem}

\begin{annexitem}{annex:p15-w02b-p15-075-stratified}{thm:p15-w02b-p15-075-stratified}{The stratified estimator under the assignment premise}
\begin{proof}
(1) Condition on the counts. By (P), the conditional law of the assignment is
uniform on the assignments with \(n_{1b}\) treated units in each stratum \(b\).
That set is a product over strata, so the conditional law is the product of
uniform laws, one per stratum. This is blocked complete randomization
(\cref{def:p15-w02b-p15-075-blocked}) with two levels, blocks equal to strata,
and cell sizes \(n_{1b},n_{0b}\geq1\) by (Q). The estimator
\(\widehat\tau_{\rm str}\) is \(\widehat\tau_{\mathcal B}\), and
\cref{thm:p15-w02b-p15-075-blocked} gives all three statements.

(2) In \cref{thm:p15-w02b-p15-075-worked}(7), (P) and (Q) hold and
\(\mathbb E[D]=4/9\neq3/2=\tau\).

(3) Under (P), every assignment with the realized counts has probability
\(1/\prod_b\binom{N_b}{n_{1b}}>0\) given the counts. Under the point mass at
\(z\), the realized assignment has probability one. The constraints that the
observed \((z,y)\) place on the table are \(Y_i(z_i)=y_i\). They do not mention
the law. So the same tables, and the same range of \(\tau\) from
\cref{thm:p15-w02b-p15-075-observational}(3), are compatible with the observed data
under either law.
\end{proof}
\end{annexitem}

\begin{annexitem}{annex:p15-w02b-p15-075-worked}{thm:p15-w02b-p15-075-worked}{Verification of the six-unit calculations}
\begin{proof}
(1) The contrasts \(2,1,1,2,0,3\) sum to \(9\), so \(\tau=3/2\). The values
\(Y_i(1)\) have mean \(5\) and squared deviations \(4,4,1,1,0,16\), with sum
\(26\), so \(S_1^2=26/5\). The values \(Y_i(0)\) have mean \(7/2\) and
squared deviations summing to \(35/2\), so \(S_0^2=7/2\). The contrasts have
squared deviations \(1/4,1/4,1/4,1/4,9/4,9/4\), with sum \(11/2\), so
\(S_w^2=11/10\). The severe stratum's contrasts \(2,1,1\) average
\(4/3\), and the others' contrasts \(2,0,3\) average \(5/3\).

(2) By \cref{thm:p15-w02b-p15-075-neyman},
\(\operatorname{Var}(\widehat\tau)=26/15+7/6-11/60=163/60\) and
\(\mathbb E\widehat V=26/15+7/6=29/10\). For the treated set \(\{1,4,5\}\),
the observed treated responses are \(3,6,5\), with mean \(14/3\) and
\(s_1^2=7/3\). The untreated responses are \(2,3,6\), with mean \(11/3\) and
\(s_0^2=13/3\). Thus \(\widehat\tau=1\) and
\(\widehat V=(7/3+13/3)/3=20/9\).

(3) Under \(H_0\) the response vector is \((3,2,3,6,5,6)\) for every
assignment, with total \(25\). For a treated set \(A\) of size three with sum
\(r\), the statistic is \(r/3-(25-r)/3=(2r-25)/3\). The twenty treated sets,
with \(r\) and the statistic, are
\[
\begin{array}{llll}
\{1,2,3\}:8,\,-3 & \{1,2,4\}:11,\,-1 & \{1,2,5\}:10,\,-\tfrac53 & \{1,2,6\}:11,\,-1\\
\{1,3,4\}:12,\,-\tfrac13 & \{1,3,5\}:11,\,-1 & \{1,3,6\}:12,\,-\tfrac13 & \{1,4,5\}:14,\,1\\
\{1,4,6\}:15,\,\tfrac53 & \{1,5,6\}:14,\,1 & \{2,3,4\}:11,\,-1 & \{2,3,5\}:10,\,-\tfrac53\\
\{2,3,6\}:11,\,-1 & \{2,4,5\}:13,\,\tfrac13 & \{2,4,6\}:14,\,1 & \{2,5,6\}:13,\,\tfrac13\\
\{3,4,5\}:14,\,1 & \{3,4,6\}:15,\,\tfrac53 & \{3,5,6\}:14,\,1 & \{4,5,6\}:17,\,3
\end{array}
\]
Counting gives the stated multiplicities. The observed value is \(1\). Eight
sets have statistic at least \(1\), giving \(8/20=2/5\), and sixteen have
absolute value at least \(1\), giving \(16/20=4/5\).

(4) Under \(H_\delta\), \(Y_i(0)=a_i\) with \(a=y-\delta z\) and
\(z\) the realized assignment. For a treated set \(A\), the imputed responses
are \(a+\delta\mathbf1_A\), and the statistic is
\(|\bar a_A-\bar a_{A^c}|\), where bars denote means. Writing \(s(A)\) for
the entry of the table in~(3) and \(k=|A\cap\{1,4,5\}|\),
\(\bar a_A-\bar a_{A^c}=s(A)-\delta(2k-3)/3\). The realized set has
\(k=3\) and its complement \(k=0\); both give \(|1-\delta|\). Complementary
sets give equal absolute values, and each pair has one member with \(k=2\).
So \(p_{H_\delta}=2/20\) unless some set with \(k=2\) satisfies
\(|s(A)-\delta/3|\geq|1-\delta|\), in which case \(p_{H_\delta}\geq4/20\).
The sets with \(k=2\) have \(s(A)\in\{-5/3,-1,-1/3,1/3,1,5/3,3\}\). For
\(s=3\), the condition \(|3-\delta/3|\geq|1-\delta|\) holds exactly for
\(-3\leq\delta\leq3\). For \(\delta\geq1\) and \(\delta\leq9\) it reads
\(3-\delta/3\geq\delta-1\). For \(\delta<1\) it reads
\(3-\delta/3\geq1-\delta\). For \(\delta>9\) it would need
\(\delta/3-3\geq\delta-1\), which fails. For \(s=-5/3\),
\(|5+\delta|/3\geq|1-\delta|\) holds exactly for \(-1/2\leq\delta\leq4\),
by the same case analysis at \(\delta=1\) and \(\delta=-5\). For
\(s\in\{-1,-1/3,1/3,1,5/3\}\) the condition reads \(|3+\delta|\geq3|1-\delta|\),
\(|1+\delta|\geq3|1-\delta|\), \(|1-\delta|\geq3|1-\delta|\),
\(|3-\delta|\geq3|1-\delta|\) or \(|5-\delta|\geq3|1-\delta|\), with solution
sets \([0,3]\), \([1/2,2]\), \(\{1\}\), \([0,3/2]\) and \([-1,2]\). These all
lie inside \([-3,4]\). The union is \([-3,3]\cup[-1/2,4]=[-3,4]\). Inside it
\(p_{H_\delta}\geq1/5>1/10\); outside it \(p_{H_\delta}=1/10\). Hence
\(C=[-3,4]\).

(5) \(B=\frac65\cdot\frac{100}4\cdot\bigl(\frac13+\frac13\bigr)=20\), and
\(\sqrt{B/\eta}=\sqrt{100}=10\).

(6) The treated responses are \(Y_i(1)=3,3,4\), with mean \(10/3\). The
untreated responses are \(Y_i(0)=4,5,6\), with mean \(5\). So \(D=-5/3\). The
treated contrasts \(2,1,1\) average \(4/3\), and
\(\sigma=2-5=-3\). With sums \(10\) and \(15\), \(n_1=n_0=3\) and \([0,10]\),
\(L=(10-15+0-30)/6=-35/6\) and \(U=(10-15+30-0)/6=25/6\).

(7) Each severe unit is treated with probability \(2/3\) and each other unit
with probability \(1/3\), and \(n_1=n_0=3\). So
\(\mathbb E\overline Y_1=\tfrac13\bigl(\tfrac23(3+3+4)+\tfrac13(6+5+9)\bigr)=40/9\)
and \(\mathbb E\overline Y_0=\tfrac13\bigl(\tfrac13(1+2+3)+\tfrac23(4+5+6)\bigr)=4\).
Hence \(\mathbb E[D]=4/9\). The stratified estimator has expectation \(\tau\)
by \cref{thm:p15-w02b-p15-075-blocked}(1). For its variance, in the severe
stratum \(S_{1}^2=1/3\), \(S_0^2=1\), \(S_w^2=1/3\), with
\(n_1=2\), \(n_0=1\), so \(V=1/6+1-1/9=19/18\). In the other stratum
\(S_1^2=13/3\), \(S_0^2=1\), \(S_w^2=7/3\), with \(n_1=1\), \(n_0=2\), so
\(V=13/3+1/2-7/9=73/18\). Then
\(\operatorname{Var}=\tfrac14\cdot\tfrac{19}{18}+\tfrac14\cdot\tfrac{73}{18}=23/18\).

(8) Here \(q_{\mathcal R}(1)=q_{\mathcal R}(0)=1/2\). By
\cref{thm:p15-w02b-p15-075-estimand-gap}(2),
\(\tau_{\mathcal U}=\tfrac15\cdot\tfrac43+\tfrac45\cdot\tfrac53=8/5\).
The composition term is
\((\tfrac15-\tfrac12)\tfrac43+(\tfrac45-\tfrac12)\tfrac53=-\tfrac25+\tfrac12=\tfrac1{10}\),
and the within-stratum term is zero by the premise. For the nested case,
\cref{thm:p15-w02b-p15-075-estimand-gap}(3) gives endpoints
\(\tfrac6{30}\cdot\tfrac32\mp\tfrac{24}{30}\cdot10=\tfrac3{10}\mp8\).
\end{proof}
\end{annexitem}
